% !TeX program = pdflatex
%=======================================================================
%  Brief de investigacion - robots de inspeccion industrial
%  Compilar:  pdflatex archivo.tex   (dos veces, para el indice)
%=======================================================================
\documentclass[11pt,a4paper]{article}

\usepackage[T1]{fontenc}
\usepackage[utf8]{inputenc}
\usepackage[spanish,es-noshorthands,es-tabla]{babel}
\usepackage{lmodern}
\usepackage[a4paper,top=2.6cm,bottom=2.6cm,left=3cm,right=2.6cm]{geometry}

\usepackage{xcolor}
\usepackage{graphicx}
\usepackage{booktabs}
\usepackage{longtable}
\usepackage{pdflscape}
\usepackage{array}
\usepackage{ragged2e}
\usepackage{enumitem}
\usepackage{amssymb}
\usepackage{titlesec}
\usepackage{fancyhdr}
\usepackage{microtype}
\usepackage{parskip}
\usepackage[hidelinks,breaklinks=true]{hyperref}
\usepackage{url}

% --- colores ----------------------------------------------------------
\definecolor{accent}{HTML}{14496B}
\definecolor{rulegray}{HTML}{C9D0D9}
\definecolor{calloutbg}{HTML}{EEF3F8}
\definecolor{inksoft}{HTML}{4A5568}

\hypersetup{colorlinks=true, linkcolor=accent, urlcolor=accent,
            citecolor=accent, pdftitle={Brief de investigación para Claude Code}}

% --- URLs largas ------------------------------------------------------
\renewcommand{\UrlFont}{\small\ttfamily}
\Urlmuskip=0mu plus 1mu
\def\UrlBreaks{\do\/\do\-\do\_\do\.\do\=\do\&\do\?\do\#}
\sloppy

% --- tablas: columnas justificadas, ancho calculado por tabla ----------
\newlength{\colw}
\newcolumntype{L}{>{\RaggedRight\arraybackslash}p{\colw}}
\setlength{\LTcapwidth}{\linewidth}
\renewcommand{\arraystretch}{1.25}

% --- listas -----------------------------------------------------------
\setlist{topsep=3pt, itemsep=1.5pt, parsep=0pt, leftmargin=1.4em}

% --- titulos ----------------------------------------------------------
\titleformat{\section}
  {\sffamily\Large\bfseries\color{accent}}{}{0pt}{}
  [\vspace{2pt}{\color{rulegray}\titlerule[0.8pt]}]
\titleformat{\subsection}{\sffamily\large\bfseries}{}{0pt}{}
\titleformat{\subsubsection}{\sffamily\normalsize\bfseries\color{inksoft}}{}{0pt}{}
\titleformat{\paragraph}[runin]{\sffamily\bfseries\small}{}{0pt}{}[.\quad]
\titlespacing*{\section}{0pt}{22pt}{9pt}
\titlespacing*{\subsection}{0pt}{16pt}{6pt}
\titlespacing*{\subsubsection}{0pt}{12pt}{4pt}

% --- encabezado / pie -------------------------------------------------
\pagestyle{fancy}
\fancyhf{}
\renewcommand{\headrulewidth}{0.4pt}
\fancyhead[L]{\sffamily\footnotesize\color{inksoft}Brief de investigación --- Robots de inspección industrial}
\fancyfoot[C]{\sffamily\footnotesize\thepage}

% --- callout ----------------------------------------------------------
\newenvironment{callout}
  {\par\medskip\noindent
   \begin{lrbox}{\calloutbox}%
   \begin{minipage}{\dimexpr\linewidth-2.2em\relax}\itshape}
  {\end{minipage}\end{lrbox}%
   \noindent{\color{accent}\vrule width 2.5pt}\hspace{0.6em}%
   \colorbox{calloutbg}{\usebox{\calloutbox}}\par\medskip}
\newsavebox{\calloutbox}

\begin{document}
\pagenumbering{roman}

%=======================================================================
% Portada
%=======================================================================
\begin{titlepage}
\thispagestyle{empty}
\setlength{\parindent}{0pt}
\raggedright
\vspace*{3.2cm}
{\sffamily\footnotesize\bfseries\color{accent}%
 \MakeUppercase{Brief de investigación}\par}
\vspace{8pt}
{\color{accent}\rule{\linewidth}{2.5pt}}
\vspace{20pt}

{\sffamily\huge\bfseries Brief de investigación para Claude Code\par}
\vspace{12pt}
{\sffamily\large\color{inksoft} Mapa exhaustivo de oportunidades de robots de inspección industrial para generación y utilities\par}
\vspace{30pt}
{\color{rulegray}\rule{0.35\linewidth}{0.6pt}\par}
\vspace{14pt}
{\normalsize \textbf{Idioma del informe final:} español\\[2pt]
\textbf{Fecha de corte inicial:} 31 de agosto de 2026\\[2pt]
\textbf{Objetivo de salida:} informe profesional HTML autocontenido, con rigor de \emph{state of the art} y análisis de negocio, comparable en disciplina al informe previo de Wedge Tightness Testing (WTT).\par}
\vfill
{\sffamily\footnotesize\color{inksoft}Documento de trabajo --- brief metodológico.
Las hipótesis y fuentes seed aquí listadas no constituyen conclusiones de
investigación.\par}
\end{titlepage}
\setcounter{page}{2}

\tableofcontents
\clearpage
\pagenumbering{arabic}

%=======================================================================
% Cuerpo
%=======================================================================

\section*{0. Tu misión}
\addcontentsline{toc}{section}{0. Tu misión}
Actuá como un equipo combinado de:

\begin{itemize}
\item ingeniero senior de robótica de inspección;
\item especialista en NDE/NDT;
\item analista de mercado industrial B2B;
\item analista de mantenimiento de generación eléctrica;
\item investigador de prior art/patentes;
\item analista de producto y estrategia.
\end{itemize}

Debés producir una investigación \textbf{profunda, verificable y escéptica} sobre qué robots de inspección industrial conviene desarrollar como próximos productos.

La pregunta central NO es:

\begin{callout}
``¿Qué robots interesantes podríamos construir?''
\end{callout}

La pregunta es:

\begin{callout}
\textbf{``¿En qué tareas de inspección/mantenimiento existe un dolor económico suficientemente grande, acceso suficientemente difícil y un mercado suficientemente abierto como para justificar el desarrollo de un robot industrial especializado y construir un negocio rentable alrededor de él?''}
\end{callout}

La referencia de calidad metodológica es un \emph{prior-art review}: identificar qué existe, qué se probó, qué falló, quién lo vende, cuánto valor crea, qué barreras quedan y dónde existe realmente un espacio defendible para un nuevo producto.

\textbf{No des por válida ninguna de las oportunidades preliminares de este brief. Verificalas desde cero y descartalas si la evidencia lo exige.}


\section*{1. Contexto del producto de referencia}
\addcontentsline{toc}{section}{1. Contexto del producto de referencia}
Ya existe un proyecto de robot RI (\emph{robotic inspection}) para inspección de generadores con rotor instalado. Usalo como \textbf{benchmark económico y tecnológico}, no como oportunidad nueva.

La lógica del negocio de referencia es:

\begin{itemize}
\item activo extremadamente caro;
\item acceso difícil;
\item una inspección convencional puede exigir gran desmontaje, personal especializado o una parada larga;
\item el robot reduce outage, desmontaje y riesgo;
\item el cliente no compra solamente ``un robot'': compra \textbf{información de condición confiable} y ahorro de mantenimiento;
\item el valor del servicio puede ser muy superior al costo del hardware;
\item la plataforma combina locomoción, cámaras, iluminación, sensado especializado, posicionamiento, tether/comunicaciones, procesamiento y generación automática de informes.
\end{itemize}

Buscamos oportunidades con una lógica económica comparable.


\subsection*{Restricciones y filosofía de producto}
\addcontentsline{toc}{subsection}{Restricciones y filosofía de producto}
Al evaluar soluciones propuestas, priorizar:

\begin{enumerate}
\item \textbf{Calidad industrial:} hardware apto para servicio profesional repetido y alta confiabilidad.
\item \textbf{Costo/beneficio:} evitar soluciones metrológicas o robóticas innecesariamente exóticas si una arquitectura commodity logra el resultado.
\item \textbf{Componentes commodity:} preferir componentes con múltiples fabricantes y cadena de suministro amplia.
\item \textbf{Industrialización:} evaluar manufacturabilidad, mantenimiento, calibración, reemplazo de consumibles y vida útil.
\item \textbf{Argentina/Latinoamérica:} considerar disponibilidad de componentes, soporte, exportación/importación y potencial de mercado regional.
\item \textbf{Precio:} buscar precios reales cuando existan. Si el OEM no publica precio, buscar distribuidores, contratos públicos, RFP/RFQ, compras gubernamentales, integradores, mercado usado o estimaciones claramente etiquetadas.
\item \textbf{FME / recuperación:} en generación eléctrica y nuclear, un robot que queda atrapado puede ser peor que no inspeccionar. La estrategia de recuperación es parte del producto.
\item \textbf{Servicio antes que hardware:} considerar seriamente Robotics-as-a-Service / Inspection-as-a-Service si la frecuencia por activo es baja pero el ticket por inspección es alto.
\end{enumerate}


\section*{2. Regla de oro: separar HECHO, DECLARACIÓN e INFERENCIA}
\addcontentsline{toc}{section}{2. Regla de oro: separar HECHO, DECLARACIÓN e INFERENCIA}
El informe debe aplicar una disciplina estricta similar a un buen \emph{engineering decision document}.

Para cada afirmación relevante, distinguir:


\subsubsection*{A. Hecho verificado}
\addcontentsline{toc}{subsubsection}{A. Hecho verificado}
Información respaldada por una fuente primaria o independiente de alta calidad.

Ejemplo:

\begin{callout}
EPRI estima ahorros \textgreater{}USD 1 M y \textless{}USD 5 M por utilización para una aplicación específica de crawler de tuberías enterradas [ref].
\end{callout}


\subsubsection*{B. Declarado por fabricante}
\addcontentsline{toc}{subsubsection}{B. Declarado por fabricante}
Una afirmación comercial del OEM/proveedor.

Ejemplo:

\begin{callout}
TesTex declara que su IAT puede inspeccionar todos los tubos conectados a un header en 3–4 turnos [ref].
\end{callout}

No convertir una afirmación del proveedor en validación independiente.


\subsubsection*{C. Inferencia de ingeniería}
\addcontentsline{toc}{subsubsection}{C. Inferencia de ingeniería}
Conclusión razonada del autor.

Debe etiquetarse explícitamente:

\begin{callout}
\textbf{Inferencia de ingeniería:} ...
\end{callout}


\subsubsection*{D. Estimación comercial}
\addcontentsline{toc}{subsubsection}{D. Estimación comercial}
Cálculo propio de TAM, ticket, ahorro o volumen.

Debe mostrar:

\begin{itemize}
\item fórmula;
\item supuestos;
\item rango;
\item fuentes de cada entrada;
\item sensibilidad.
\end{itemize}

Nunca presentar una estimación propia como dato publicado.


\section*{3. Jerarquía de evidencia}
\addcontentsline{toc}{section}{3. Jerarquía de evidencia}
Usar y mostrar una clasificación consistente.


\subsection*{Nivel A — máxima prioridad}
\addcontentsline{toc}{subsection}{Nivel A — máxima prioridad}
\begin{itemize}
\item EPRI
\item DOE / NRC / IAEA / EIA / FERC / NERC
\item IEEE / IEC / ASME / API / NACE/AMPP
\item CIGRE
\item documentos oficiales de utilities
\item papers revisados por pares
\item patentes originales
\item documentación técnica oficial del fabricante
\item licitaciones/RFP/contratos oficiales
\end{itemize}


\subsection*{Nivel B}
\addcontentsline{toc}{subsection}{Nivel B}
\begin{itemize}
\item proceedings de congresos
\item tesis
\item informes institucionales
\item presentaciones técnicas EPRI/OEM
\item casos técnicos de utilities
\item publicaciones de asociaciones industriales serias (HRSG Forum/CCJ, etc.)
\end{itemize}


\subsection*{Nivel C}
\addcontentsline{toc}{subsection}{Nivel C}
\begin{itemize}
\item páginas de proveedores de servicios
\item casos comerciales
\item notas técnicas secundarias
\item prensa especializada
\end{itemize}


\subsection*{Nivel D}
\addcontentsline{toc}{subsection}{Nivel D}
\begin{itemize}
\item agregadores de contratos
\item blogs
\item ResearchGate/Scribd
\item notas de marketing de terceros
\end{itemize}

Nivel D solamente debe servir para \textbf{descubrir la fuente primaria}. No usarlo como evidencia final si puede localizarse el original.


\section*{4. Reglas de verificación obligatorias}
\addcontentsline{toc}{section}{4. Reglas de verificación obligatorias}
\begin{enumerate}
\item \textbf{Una patente demuestra divulgación, no funcionamiento comercial.}
\begin{itemize}
\item verificar familia;
\item titular;
\item prioridad;
\item legal status;
\item expiración/abandono cuando sea relevante.
\end{itemize}
\item \textbf{Un prototipo universitario no es un producto.}
\item \textbf{Un producto anunciado no equivale a despliegue de campo.}
Buscar:
\begin{itemize}
\item clientes;
\item sitios;
\item fechas;
\item cantidad de inspecciones;
\item evidencia de repetición.
\end{itemize}
\item \textbf{Una reducción de costos publicada por el fabricante debe identificarse como tal.}
\item \textbf{No usar market reports opacos} (``market size USD X billion'') como base principal.
\item Si no existe precio:
\begin{itemize}
\item escribir ``precio no publicado'';
\item luego buscar contratos/RFPs/distribuidores;
\item no inventar.
\end{itemize}
\item Si no hay evidencia de competencia:
\begin{itemize}
\item escribir ``no se identificó competencia después de las búsquedas X/Y/Z'';
\item no escribir ``no existe competencia''.
\end{itemize}
\item Buscar productos \textbf{actuales, discontinuados, adquiridos y fallidos}.
\item Buscar en inglés y, cuando corresponda, también:
\begin{itemize}
\item alemán;
\item japonés;
\item chino;
\item coreano;
\item francés;
\item español/portugués.
\end{itemize}

Empresas importantes pueden no estar bien indexadas en búsquedas en inglés.
\item Para cada oportunidad hacer al menos:
\begin{itemize}
\item búsquedas EPRI por activo;
\item búsquedas EPRI por ``robot/robotic/remote inspection/automation'';
\item Google/web general;
\item Google Patents/WIPO/Espacenet;
\item OEMs;
\item utilities;
\item papers IEEE/ASME/Elsevier/Springer;
\item RFP/RFQ/contratos;
\item videos técnicos sólo como evidencia suplementaria.
\end{itemize}
\end{enumerate}


\section*{5. Error que este estudio debe evitar}
\addcontentsline{toc}{section}{5. Error que este estudio debe evitar}
No queremos descubrir dentro de dos años que:

\begin{itemize}
\item EPRI ya investigó exactamente el problema;
\item un OEM ya patentó la arquitectura;
\item existe un producto comercial establecido;
\item el cliente ya compra ese servicio;
\item hubo un programa de desarrollo que fracasó por una razón física conocida.
\end{itemize}

Por eso, \textbf{para cada candidato debés hacer una búsqueda histórica además de una búsqueda actual}.

Construí una línea temporal:

\begin{callout}
problema identificado $\rightarrow$ primer prototipo $\rightarrow$ investigación EPRI $\rightarrow$ patentes $\rightarrow$ primeros despliegues $\rightarrow$ productos comerciales $\rightarrow$ estado 2026.
\end{callout}

Las fuentes negativas son especialmente valiosas:

\begin{itemize}
\item productos discontinuados;
\item programas abandonados;
\item papers con resultados insuficientes;
\item robots que quedaron atascados;
\item herramientas que no entraron;
\item FME;
\item fallas de tether;
\item NDE que no logró sensibilidad;
\item costos que destruyeron el business case.
\end{itemize}


\section*{6. Fuentes EPRI que DEBÉS recorrer antes de concluir}
\addcontentsline{toc}{section}{6. Fuentes EPRI que DEBÉS recorrer antes de concluir}
No limitarte a estas. Son un paquete inicial verificado.


\subsection*{6.1 Mapa cross-sector de robótica EPRI}
\addcontentsline{toc}{subsection}{6.1 Mapa cross-sector de robótica EPRI}

\subsubsection*{EPRI 3002023899}
\addcontentsline{toc}{subsubsection}{EPRI 3002023899}
\textbf{Robotic Process Automation for Nuclear Power Plants: Evaluation of Near-Term Opportunities} \\
Junio 2022 \\
\url{https://restservice.epri.com/publicdownload/000000003002023899/0/Product}

Usarlo para:

\begin{itemize}
\item lista de tareas de alto valor;
\item horas-hombre;
\item gaps tecnológicos;
\item diferencia entre teleoperación, autonomía e inspección+mantenimiento;
\item criterio de selección de tareas.
\end{itemize}

Este documento identifica, entre otras, oportunidades en:

\begin{itemize}
\item buried service water piping;
\item component supports;
\item cooling-water intake/discharge structures;
\item isophase bus duct;
\item 3D radiation mapping;
\item containment leak detection;
\item reactor head studs/stud holes;
\item water storage tanks;
\item dry fuel casks;
\item operator rounds;
\item fire protection;
\item steam-generator tubing.
\end{itemize}

\textbf{No asumir que todas son buenos productos nuevos. Varias ya son mercados maduros.}


\subsubsection*{EPRI 3002025693}
\addcontentsline{toc}{subsubsection}{EPRI 3002025693}
\textbf{Program on Technology Innovation: Landscape of Automation in Nuclear Power Plants} \\
\url{https://restservice.epri.com/publicdownload/000000003002025693/0/Product}

Buscar especialmente la conclusión de negocio sobre:

\begin{itemize}
\item robots custom para tareas de alto valor y baja frecuencia;
\item uso frecuente de vendors externos;
\item costo/tiempo de desarrollar plataformas custom.
\end{itemize}

Esto puede respaldar un modelo de \textbf{servicio especializado}.


\subsubsection*{EPRI Unmanned Mobile Technologies Collaboration Group}
\addcontentsline{toc}{subsubsection}{EPRI Unmanned Mobile Technologies Collaboration Group}
Product Index: \\
\url{https://transmission.epri.com/p37_substations/robotics/mobiletechgrp/products/}

Webcasts / actividad actual: \\
\url{https://transmission.epri.com/p37_substations/robotics/mobiletechgrp/webcasts/}

Recorrer \textbf{todo el índice}, no sólo los títulos que parecen obvios.


\subsubsection*{EPRI Tech Portal}
\addcontentsline{toc}{subsubsection}{EPRI Tech Portal}
\url{https://techportal.epri.com/}

Usarlo como herramienta de descubrimiento; luego ir a fuentes primarias.


\subsubsection*{EPRI Generation reports / Digital Transformation}
\addcontentsline{toc}{subsubsection}{EPRI Generation reports / Digital Transformation}
\url{https://dx-wiki.epri.com/Generation_Reports}

Buscar explícitamente:

\begin{itemize}
\item robots;
\item inspection;
\item remote inspection;
\item crawlers;
\item vine robots;
\item confined spaces;
\item NDE/NDT;
\item autonomous inspection.
\end{itemize}


\section*{7. Candidatos que deben investigarse EN PROFUNDIDAD}
\addcontentsline{toc}{section}{7. Candidatos que deben investigarse EN PROFUNDIDAD}
La lista siguiente es un punto de partida, no una conclusión.


\subsection*{CANDIDATO A — Robot para Isolated Phase Bus / Isophase Bus (IPB)}
\addcontentsline{toc}{subsection}{CANDIDATO A — Robot para Isolated Phase Bus / Isophase Bus (IPB)}

\subsubsection*{Hipótesis}
\addcontentsline{toc}{subsubsection}{Hipótesis}
Puede ser la mejor extensión inmediata del RI por:

\begin{itemize}
\item mismo tipo de cliente;
\item inspección durante outages;
\item espacio confinado;
\item reducción de scaffolding/man-entry;
\item alta reutilización de cámaras, tether, iluminación, software de mapping y arquitectura de crawler.
\end{itemize}


\subsubsection*{Fuentes EPRI iniciales}
\addcontentsline{toc}{subsubsection}{Fuentes EPRI iniciales}

\paragraph*{MTA-MA-029}
\textbf{Reduce Inspection Cost for Isophase Bus Duct Using Robotic Crawler} \\
\url{https://nuclearplantmod.epri.com/MTA-MA-029}

Datos verificados que debés revisar y citar en contexto:

\begin{itemize}
\item experiencia reportada de una utility con costo de vendor de aproximadamente \textbf{USD 100.000 por implementación};
\item inspección en \textbf{menos de un día};
\item EPRI reporta reducción aproximada de \textbf{10$\times$} frente a la inspección manual en un caso;
\item payback inmediato según SWEEP;
\item tecnología comercialmente disponible;
\item riesgos de crawler/tool failure;
\item una gran parte del ahorro proviene de evitar scaffolding y acceso humano.
\end{itemize}

No generalizar el USD 100k a todo el mercado sin evidencia.


\paragraph*{EPRI 1015057}
\textbf{Isolated Phase Bus Maintenance Guide} \\
\url{https://restservice.epri.com/publicdownload/000000000001015057/0/Product}

Buscar:

\begin{itemize}
\item frecuencia típica de inspección;
\item defectos;
\item accesos;
\item figuras de robotic walkers;
\item mantenimiento recomendado;
\item mediciones eléctricas.
\end{itemize}


\paragraph*{EPRI 3002023899}
Investigar la idea de evolucionar desde ``cámara crawler'' a:

\begin{itemize}
\item limpieza;
\item transporte/remoción de FOD;
\item verificación de torque;
\item resistance testing;
\item manipulación.
\end{itemize}


\subsubsection*{Competidores/soluciones a investigar}
\addcontentsline{toc}{subsubsection}{Competidores/soluciones a investigar}
No limitarse a:

\begin{itemize}
\item Trident NGS — \url{https://www.tridentngs.com/}
\item Electrical Builders Inc.
\item MidStates Energy
\item nVent
\item GE Vernova / Siemens Energy / Hitachi Energy, si ofrecen servicios
\item empresas nucleares especializadas
\item sistemas coreanos/japoneses/chinos
\item trabajos de Beihang y otras universidades.
\end{itemize}

Buscar:

\begin{itemize}
\item nombre exacto del crawler;
\item fabricante;
\item locomoción;
\item dimensiones;
\item tether/wireless;
\item cómo pasa aisladores/supports;
\item recuperación;
\item sensores;
\item limpieza;
\item si mide torque o resistencia;
\item precio de servicio;
\item clientes;
\item frecuencia de inspección.
\end{itemize}


\subsubsection*{Pregunta estratégica}
\addcontentsline{toc}{subsubsection}{Pregunta estratégica}
\textbf{¿Existe espacio para un sistema ``inspection + maintenance'' y no solamente otro crawler con cámara?}


\subsection*{CANDIDATO B — Robot para HRSG: headers, tube-to-header welds y tubos}
\addcontentsline{toc}{subsection}{CANDIDATO B — Robot para HRSG: headers, tube-to-header welds y tubos}

\subsubsection*{Fuente EPRI fundamental}
\addcontentsline{toc}{subsubsection}{Fuente EPRI fundamental}

\paragraph*{EPRI 1017635}
\textbf{Study for Snake Robot Technology for Inspection of Headers and Tubes in Heat Recovery Steam Generators} \\
Final Report, 2009 \\
\url{https://restservice.epri.com/publicdownload/000000000001017635/0/Product}

Leer COMPLETO.

Extraer:

\begin{itemize}
\item por qué la inspección convencional no llega a tubos interiores;
\item geometrías;
\item diámetros;
\item longitudes;
\item tipos de HRSG;
\item puntos de acceso;
\item modalidades de inspección;
\item qué conceptos de snake robot analizaron;
\item qué demostraron;
\item qué no resolvieron;
\item por qué no se comercializó directamente ese concepto;
\item NDE propuesto.
\end{itemize}


\subsubsection*{Competidor crítico: TesTex}
\addcontentsline{toc}{subsubsection}{Competidor crítico: TesTex}

\paragraph*{TesTex HRSG tools}
\url{https://testex-ndt.com/services/hrsg/}


\paragraph*{TesTex IAT}
\url{https://testex-ndt.com/products/iat-internal-access-tool-for-hrsg-inspections/}

Información declarada por TesTex a verificar:

\begin{itemize}
\item \textbf{Internal Access Tool (IAT)} es un crawler colocado dentro del header;
\item puede entrar por end cap o versión modular por un agujero de aproximadamente 8'';
\item empuja una sonda RFET + cámara por los tubos;
\item remote control;
\item puede inspeccionar todos los tubos de un header en 3–4 shifts;
\item tiene despliegues con múltiples utilities;
\item el sitio declara 100\% tube coverage dentro de las condiciones de aplicación.
\end{itemize}


\paragraph*{TesTex ``Claw''}
\url{https://testex-ndt.com/products/claw-bfet-inspection-tool/}

Investigar:

\begin{itemize}
\item BFET;
\item cobertura real;
\item 200 welds/shift declarado;
\item diámetros;
\item accesibilidad.
\end{itemize}


\subsubsection*{Investigación adicional obligatoria}
\addcontentsline{toc}{subsubsection}{Investigación adicional obligatoria}
Buscar:

\begin{itemize}
\item patents TesTex;
\item EPRI/TesTex joint development;
\item Combined Cycle Journal / HRSG Forum;
\item GE Vernova;
\item Siemens Energy;
\item Nooter/Eriksen;
\item Vogt;
\item Alstom/Altrad;
\item EDF;
\item MISTRAS;
\item Eddyfi;
\item Waygate/Baker Hughes;
\item Olympus/Evident;
\item robots académicos para boiler tubes;
\item tube crawlers magnéticos;
\item snake robots;
\item continuum robots.
\end{itemize}


\subsubsection*{Pregunta estratégica}
\addcontentsline{toc}{subsubsection}{Pregunta estratégica}
No asumir mercado vacío. Determinar:

\textbf{¿Qué parte del problema NO resuelve hoy TesTex?}

Posibles hipótesis a validar:

\begin{itemize}
\item menor access hole;
\item no retirar end cap;
\item más tipos de headers;
\item automatic tube indexing;
\item mayor alcance;
\item mejor localization;
\item un solo despliegue para weld + tube wall;
\item NDE cuantitativo superior;
\item mejor autonomía;
\item depósitos/bloqueos;
\item inspección + limpieza/reparación.
\end{itemize}

Si no aparece un gap defendible, bajar el ranking.


\subsection*{CANDIDATO C — Robot para underground transmission vaults / HV cable manholes}
\addcontentsline{toc}{subsection}{CANDIDATO C — Robot para underground transmission vaults / HV cable manholes}

\subsubsection*{EPRI histórico + actual}
\addcontentsline{toc}{subsubsection}{EPRI histórico + actual}

\paragraph*{EPRI 3002000878}
\textbf{Underground Transmission Vault Inspection Using Robotic Techniques} \\
2013 \\
Localizar en EPRI y descargar si está público.


\paragraph*{EPRI 3002032834}
\textbf{Underground Transmission Vault Inspection Using Robotic Techniques — 2025 Update}

Página de research updates:
\url{https://transmission.epri.com/p36_underground/public/p36001_design_construction/research_updates/}

Hecho crítico:
EPRI continúa este trabajo en \textbf{2026 y 2027} y declara que:

\begin{itemize}
\item se busca mejorar seguridad de trabajadores;
\item reducir requisitos de circuit outage;
\item conceptos/prototipos robóticos fueron demostrados en laboratorios EPRI y sitios de utilities.
\end{itemize}


\paragraph*{Applications}
\url{https://transmission.epri.com/p36_underground/public/p36_applications/}

La página enumera:

\begin{callout}
Robotic Inspections of Underground Vaults — Prototype Development, Laboratory Tests, Demo in Planning.
\end{callout}


\paragraph*{Contactos EPRI actuales}
\url{https://transmission.epri.com/p36_underground/leads/}

Verificar antes de usar:

\begin{itemize}
\item David Kummer — Robotic Inspection;
\item Tom Zhao — Underground Transmission Program Manager.
\end{itemize}


\subsubsection*{Proyecto complementario actual}
\addcontentsline{toc}{subsubsection}{Proyecto complementario actual}
Buscar ``Extruded Cable Manhole Inspections'' en los Collaborative Supplemental Projects 2025/2026 de EPRI.

Investigar la arquitectura:

\begin{itemize}
\item inspección desde superficie;
\item energized cable;
\item IR;
\item cámara;
\item gas sensing;
\item mechanical arm;
\item PD/acoustic/ultrasonic si aplica.
\end{itemize}


\subsubsection*{Competencia}
\addcontentsline{toc}{subsubsection}{Competencia}
Buscar en profundidad:

\begin{itemize}
\item Osmose;
\item Con Edison;
\item National Grid;
\item SCE;
\item PG\&E;
\item NYPA;
\item Duke;
\item Dominion;
\item Entergy;
\item Hydro-Québec;
\item empresas de sewer crawlers adaptados;
\item drones confined-space;
\item Boston Dynamics/ANYbotics en vaults;
\item integradores.
\end{itemize}


\subsubsection*{Pregunta estratégica}
\addcontentsline{toc}{subsubsection}{Pregunta estratégica}
Este candidato puede estar \textbf{menos maduro} que IPB o buried-pipe NDE.

Determinar:

\begin{itemize}
\item cantidad de vaults;
\item frecuencia de inspección;
\item necesidad de de-energización;
\item costo actual;
\item safety/confined-space burden;
\item si puede venderse como servicio recurrente;
\item ticket razonable;
\item requirements de aislamiento eléctrico y gas.
\end{itemize}


\subsection*{CANDIDATO D — Crawler NDE para piping enterrado / inaccesible}
\addcontentsline{toc}{subsection}{CANDIDATO D — Crawler NDE para piping enterrado / inaccesible}

\subsubsection*{EPRI MTA-MA-017}
\addcontentsline{toc}{subsubsection}{EPRI MTA-MA-017}
\url{https://nuclearplantmod.epri.com/MTA-MA-017}

Datos EPRI iniciales:

\begin{itemize}
\item inspección desde un único acceso;
\item puede evitar excavación;
\item EMAT permite determinados exámenes volumétricos sin remover coating;
\item EPRI SWEEP:
\begin{itemize}
\item implementation cost \textless{} USD 1 M;
\item ahorro \textgreater{} USD 1 M y \textless{} USD 5 M por utilización;
\item payback inmediato/\textless{}1 año;
\item tecnología ya comercial en nuclear.
\end{itemize}
\end{itemize}


\subsubsection*{EPRI 1025272}
\addcontentsline{toc}{subsubsection}{EPRI 1025272}
\textbf{Compilation of Lessons Learned on Buried and Underground Piping} \\
\url{https://restservice.epri.com/publicdownload/000000000001025272/0/Product}

Buscar:

\begin{itemize}
\item por qué conventional ILI no funciona bien en nuclear;
\item falta de launcher/receiver;
\item elbows;
\item tees;
\item verticals;
\item coatings;
\item diámetros;
\item water-filled/dry;
\item FME/recovery.
\end{itemize}


\subsubsection*{Competidor crítico: Diakont}
\addcontentsline{toc}{subsubsection}{Competidor crítico: Diakont}
\url{https://diakont.com/case-studies/nuclear-solutions/first-buried-piping-in-line-inspection/}

También buscar:

\begin{itemize}
\item Structural Integrity Associates;
\item Westinghouse;
\item Framatome;
\item Eddyfi/Inuktun;
\item Waygate;
\item Gecko Robotics;
\item MISTRAS;
\item ROSEN;
\item NDT Global;
\item Eddyfi VersaTrax;
\item pipe inspection service companies.
\end{itemize}


\subsubsection*{Pregunta estratégica}
\addcontentsline{toc}{subsubsection}{Pregunta estratégica}
EPRI ya lo clasifica como comercialmente implementado.

Por lo tanto:
\textbf{no recomendar un crawler genérico de tuberías.}

Sólo recomendar si se identifica white space concreto:

\begin{itemize}
\item small diameter;
\item 1D elbow;
\item tees;
\item reducers;
\item vertical;
\item wet/dry transition;
\item protected/coated pipe;
\item inspection + repair;
\item deployability without cutting;
\item smaller access;
\item low-cost service outside nuclear.
\end{itemize}


\subsection*{CANDIDATO E — Robot para dry spent-fuel storage casks}
\addcontentsline{toc}{subsection}{CANDIDATO E — Robot para dry spent-fuel storage casks}

\subsubsection*{EPRI}
\addcontentsline{toc}{subsubsection}{EPRI}
Buscar:
\textbf{3002008234 — Dry Canister Storage System Inspection and Robotic Delivery System Development}

EPRI Journal:
\url{https://eprijournal.com/wall-climbing-robots-inspect-nuclear-storage-casks/}

Datos iniciales:

\begin{itemize}
\item gaps pequeños;
\item cámara + ECA + EMAT + radiation/temp en ciertas configuraciones;
\item múltiples field tests;
\item EPRI describe reducción de costos muy alta frente a métodos con heavy lift/movimiento de canister.
\end{itemize}


\subsubsection*{Competidor crítico}
\addcontentsline{toc}{subsubsection}{Competidor crítico}
\textbf{Robotic Technologies of Tennessee (RTT)} \\
\url{https://www.robotictechtn.com/}

El fabricante declara que nuclear sites usan sus plataformas para Dry Cask Storage inspection.

Buscar:

\begin{itemize}
\item robots exactos;
\item EPRI/DOE programs;
\item SONGS;
\item Holtec;
\item Orano;
\item NAC;
\item DOE;
\item NRC requirements;
\item Penn State/PRINSE;
\item radiation qualification;
\item patent landscape.
\end{itemize}


\subsubsection*{Pregunta estratégica}
\addcontentsline{toc}{subsubsection}{Pregunta estratégica}
Alto valor, pero:

\begin{itemize}
\item mercado nuclear limitado;
\item qualification;
\item FME;
\item radiation;
\item incumbent.
\end{itemize}

Compararlo rigurosamente contra oportunidades menos reguladas.


\subsection*{CANDIDATO F — ROV/AUV/crawler para penstocks y túneles hidro}
\addcontentsline{toc}{subsection}{CANDIDATO F — ROV/AUV/crawler para penstocks y túneles hidro}

\subsubsection*{EPRI TR-113584-V7 / 1007576}
\addcontentsline{toc}{subsubsection}{EPRI TR-113584-V7 / 1007576}
\textbf{ROV Technology: Applications and Advancements at Hydro Facilities} \\
\url{https://restservice.epri.com/publicdownload/000000000001007576/0/Product}

Extraer:

\begin{itemize}
\item tunnel/penstock;
\item visual;
\item sonar;
\item UT thickness;
\item tether limitations;
\item alcance;
\item recuperación;
\item case studies;
\item ahorro de outage/dewatering.
\end{itemize}

Buscar además:
\textbf{3002011682 — Autonomous Underwater Vehicles for Tunnel and Penstock Inspection}

Y proyectos actuales del EPRI Unmanned Mobile Technologies Collaboration Group.


\subsubsection*{Competencia}
\addcontentsline{toc}{subsubsection}{Competencia}
Investigar:

\begin{itemize}
\item Deep Trekker;
\item VideoRay;
\item Eddyfi/Inuktun;
\item Deep Ocean Engineering;
\item Saab Seaeye;
\item Blue Robotics integrators;
\item utility-owned custom ROVs;
\item Hydro-Québec;
\item Voith;
\item Andritz;
\item GE Vernova Hydro;
\item specialized tunnel inspection vendors.
\end{itemize}


\subsubsection*{Pregunta estratégica}
\addcontentsline{toc}{subsubsection}{Pregunta estratégica}
ROV visual puede ser maduro.

El posible gap puede estar en:

\begin{itemize}
\item navegación autónoma en km;
\item localization sin GPS;
\item wall thickness mapping;
\item close-contact NDE;
\item recovery guaranteed;
\item turbidity;
\item high flow;
\item long tether.
\end{itemize}


\subsection*{CANDIDATO G — Vine / everting robot para equipos de planta}
\addcontentsline{toc}{subsection}{CANDIDATO G — Vine / everting robot para equipos de planta}

\subsubsection*{EPRI}
\addcontentsline{toc}{subsubsection}{EPRI}
Buscar y verificar:
\textbf{3002032954 — Everting Vine Robots for Plant Equipment Inspection – Technology Review and Feasibility Study}

Catálogo:
\url{https://dx-wiki.epri.com/Generation_Reports}

No asumir resultados no públicos.


\subsubsection*{Fuentes externas iniciales}
\addcontentsline{toc}{subsubsection}{Fuentes externas iniciales}
\begin{itemize}
\item IEEE Robotics \& Automation Magazine — large-scale vine robots for industrial inspection.
\item UCSB/Bechtel capstone.
\item Stanford vine robots.
\item Trellis Robotics.
\item IvySpec.
\end{itemize}


\subsubsection*{Pregunta estratégica}
\addcontentsline{toc}{subsubsection}{Pregunta estratégica}
``Vine robot'' es una tecnología, NO un mercado.

Identificar un activo concreto donde:

\begin{itemize}
\item crawler no pasa;
\item borescope no alcanza;
\item desmontaje es costoso;
\item el vine pueda llevar NDE útil.
\end{itemize}

Ejemplos a evaluar:

\begin{itemize}
\item ductos de aire/gases;
\item HRSG;
\item piping complejo;
\item turbine/boiler spaces;
\item cable ducts;
\item confined equipment.
\end{itemize}


\section*{8. Mercados que deben analizarse como ``posiblemente demasiado maduros''}
\addcontentsline{toc}{section}{8. Mercados que deben analizarse como ``posiblemente demasiado maduros''}
No descartarlos sin investigar, pero exigir evidencia de white space.


\subsection*{8.1 Transformer swimming robot}
\addcontentsline{toc}{subsection}{8.1 Transformer swimming robot}
Benchmark:
\textbf{ABB TXplore}

Buscar fuentes ABB oficiales y EPRI.

ABB ha publicado:

\begin{itemize}
\item inspección en menos de un día;
\item eliminación de tareas de drenaje/procesamiento de aceite;
\item reducción de costos de hasta aproximadamente 50\% en comunicaciones de fabricante/casos.
\end{itemize}

Pregunta:
¿hay algo que TXplore no haga y tenga mercado independiente de ABB?


\subsection*{8.2 Substation patrol robots}
\addcontentsline{toc}{subsection}{8.2 Substation patrol robots}
EPRI trabaja con plataformas comerciales, incluyendo Spot en ciertos proyectos.

Buscar:
\url{https://transmission.epri.com/p37_substations/public/robotics/substation_robot_application/}

Hipótesis:
el valor puede estar más en \textbf{sensor payload + analytics} que en fabricar otro cuadrúpedo/UGV.


\subsection*{8.3 Transmission-line crawling robots}
\addcontentsline{toc}{subsection}{8.3 Transmission-line crawling robots}
Investigar:

\begin{itemize}
\item EPRI Ti;
\item Hydro-Québec LineScout / LineROVer;
\item HiBot Expliner;
\item drones modernos.
\end{itemize}

Determinar si los drones hicieron económicamente obsoleto el crawler para visual routine inspection.


\subsection*{8.4 Tank inspection robots}
\addcontentsline{toc}{subsection}{8.4 Tank inspection robots}
Investigar:

\begin{itemize}
\item Square Robot;
\item Newton Labs;
\item Eddyfi;
\item ROVs.
\end{itemize}

Probablemente mercado maduro.


\subsection*{8.5 Reactor internals / steam-generator tube robotics}
\addcontentsline{toc}{subsection}{8.5 Reactor internals / steam-generator tube robotics}
Investigar:

\begin{itemize}
\item Framatome;
\item Westinghouse;
\item AREVA SUSI;
\item KHNP/KEPRI;
\item SG manipulators.
\end{itemize}

Es un mercado de altísimo valor, pero probablemente muy maduro y regulado.


\section*{9. Buscar oportunidades QUE NO FIGURAN en esta lista}
\addcontentsline{toc}{section}{9. Buscar oportunidades QUE NO FIGURAN en esta lista}
Esto es obligatorio.

Hacé un barrido independiente por:


\subsection*{EPRI}
\addcontentsline{toc}{subsection}{EPRI}
\begin{itemize}
\item Generation
\item Nuclear
\item Transmission
\item Distribution
\item Hydropower
\item Fossil/combined-cycle
\item Wind
\item Solar
\item substations
\item underground transmission
\item plant modernization
\item NDE
\item digital transformation
\item robotic collaboration groups
\end{itemize}


\subsection*{Problemas físicos}
\addcontentsline{toc}{subsection}{Problemas físicos}
Buscar combinaciones de:

\begin{itemize}
\item confined space;
\item inaccessible inspection;
\item rotor-in-place;
\item no disassembly;
\item no scaffolding;
\item no excavation;
\item no draining;
\item no diving;
\item no de-energization;
\item in-situ;
\item online inspection;
\item internal inspection;
\item robotic NDE;
\item remote NDE;
\item autonomous inspection;
\item inspection + repair;
\item inspection + cleaning;
\item FOD retrieval.
\end{itemize}


\subsection*{Otros activos a evaluar}
\addcontentsline{toc}{subsection}{Otros activos a evaluar}
No asumir que son buenos, pero verificarlos:

\begin{itemize}
\item air-cooled condensers;
\item condensers;
\item cooling towers;
\item boiler waterwalls;
\item boiler headers;
\item furnace internals;
\item gas turbine exhaust/combustion areas;
\item steam turbine internals;
\item wind turbine blades/towers;
\item offshore wind monopiles;
\item dams;
\item spillways;
\item intake gates;
\item large valves;
\item heat exchangers;
\item stacks/chimneys;
\item industrial furnaces;
\item petrochemical vessels/pipes;
\item mining assets si la tecnología es transferible.
\end{itemize}

Queremos descubrir \textbf{candidatos nuevos} si la evidencia los hace mejores.


\section*{10. Para CADA robot/candidato usar una ficha fija}
\addcontentsline{toc}{section}{10. Para CADA robot/candidato usar una ficha fija}
Esto es obligatorio para poder comparar.


\subsection*{[Nombre del candidato]}
\addcontentsline{toc}{subsection}{[Nombre del candidato]}

\subsubsection*{1. Activo y problema}
\addcontentsline{toc}{subsubsection}{1. Activo y problema}
\begin{itemize}
\item qué se inspecciona;
\item mecanismo de falla;
\item por qué importa.
\end{itemize}


\subsubsection*{2. Cómo se inspecciona hoy}
\addcontentsline{toc}{subsubsection}{2. Cómo se inspecciona hoy}
\begin{itemize}
\item rotor-out/disassembly/scaffolding/diver/human entry/etc.;
\item frecuencia;
\item duración;
\item personal;
\item outage;
\item riesgos.
\end{itemize}


\subsubsection*{3. Pain económico verificable}
\addcontentsline{toc}{subsubsection}{3. Pain económico verificable}
\begin{itemize}
\item costo directo;
\item costo de outage;
\item costo de acceso;
\item horas-hombre;
\item pérdida por falla;
\item ejemplos reales.
\end{itemize}

Separar:

\begin{itemize}
\item datos publicados;
\item estimaciones.
\end{itemize}


\subsubsection*{4. Propuesta de robot}
\addcontentsline{toc}{subsubsection}{4. Propuesta de robot}
Definir funcionalmente:

\begin{itemize}
\item locomoción;
\item dimensiones aproximadas;
\item sensores;
\item NDE;
\item tether/wireless;
\item autonomía;
\item mapping/localization;
\item manipulación;
\item recovery.
\end{itemize}

No diseñar con detalle todavía; definir requirements.


\subsubsection*{5. Qué valor crea}
\addcontentsline{toc}{subsubsection}{5. Qué valor crea}
\begin{itemize}
\item evita qué operación;
\item reduce cuánto tiempo;
\item aumenta cobertura;
\item mejora confiabilidad;
\item disminuye riesgo.
\end{itemize}


\subsubsection*{6. Robots/productos existentes}
\addcontentsline{toc}{subsubsection}{6. Robots/productos existentes}
Tabla obligatoria:

\begin{landscape}
\begingroup\footnotesize
\setlength{\tabcolsep}{2pt}
\setlength{\emergencystretch}{3em}
\hyphenpenalty=50 \exhyphenpenalty=50
\setlength{\colw}{\dimexpr(\linewidth-20\tabcolsep)/10\relax}
\begin{longtable}{@{}L L L L L L L L L L@{}}
\toprule
\textbf{Producto} & \textbf{Fabricante} & \textbf{Año} & \textbf{Estado 2026} & \textbf{Locomoción} & \textbf{Payload\slash NDE} & \textbf{Dimensiones} & \textbf{Despliegues} & \textbf{Precio\slash servicio} & \textbf{Fuente} \\
\midrule
\endfirsthead
\toprule
\textbf{Producto} & \textbf{Fabricante} & \textbf{Año} & \textbf{Estado 2026} & \textbf{Locomoción} & \textbf{Payload\slash NDE} & \textbf{Dimensiones} & \textbf{Despliegues} & \textbf{Precio\slash servicio} & \textbf{Fuente} \\
\midrule
\endhead
\bottomrule
\endfoot
\multicolumn{10}{@{}l@{}}{\itshape Tabla a completar durante la investigación.} \\
\end{longtable}
\endgroup
\end{landscape}

No dejar fuera:

\begin{itemize}
\item products;
\item service-only systems;
\item prototypes;
\item discontinued;
\item patents.
\end{itemize}


\subsubsection*{7. Cómo resolvieron los existentes el problema}
\addcontentsline{toc}{subsubsection}{7. Cómo resolvieron los existentes el problema}
Analizar ingeniería:

\begin{itemize}
\item adhesion;
\item wheels/tracks/magnetic;
\item snake/continuum;
\item tether;
\item power;
\item sensor coupling;
\item NDE contact;
\item crossing obstacles;
\item alignment;
\item recovery;
\item FME.
\end{itemize}


\subsubsection*{8. Qué NO resuelven}
\addcontentsline{toc}{subsubsection}{8. Qué NO resuelven}
Identificar whitespace real.


\subsubsection*{9. Mercado y madurez}
\addcontentsline{toc}{subsubsection}{9. Mercado y madurez}
Usar una escala explicada:

\begin{itemize}
\item \textbf{M0:} sólo papers/conceptos
\item \textbf{M1:} prototipos de laboratorio
\item \textbf{M2:} demostración industrial aislada
\item \textbf{M3:} servicios comerciales iniciales
\item \textbf{M4:} múltiples proveedores/despliegues
\item \textbf{M5:} mercado maduro/commodity
\end{itemize}

No confundir TRL con madurez comercial.


\subsubsection*{10. Barreras técnicas}
\addcontentsline{toc}{subsubsection}{10. Barreras técnicas}
\begin{itemize}
\item geometría;
\item temperatura;
\item radiación;
\item EMI;
\item fluidos;
\item suciedad;
\item corrosión;
\item tether;
\item adhesion;
\item battery;
\item sensing;
\item NDE calibration.
\end{itemize}


\subsubsection*{11. Riesgo de FME/atasco y recuperación}
\addcontentsline{toc}{subsubsection}{11. Riesgo de FME/atasco y recuperación}
Proponer filosofía fail-safe.


\subsubsection*{12. Regulación/calificación}
\addcontentsline{toc}{subsubsection}{12. Regulación/calificación}
\begin{itemize}
\item nuclear;
\item NDE qualification;
\item electrical HV;
\item confined-space;
\item pressure boundary;
\item ASME/API/etc.
\end{itemize}


\subsubsection*{13. Patentes/IP}
\addcontentsline{toc}{subsubsection}{13. Patentes/IP}
\begin{itemize}
\item principales familias;
\item propietarios;
\item estado legal;
\item riesgo de bloqueo;
\item posibles áreas libres.
\end{itemize}


\subsubsection*{14. Desarrollo estimado}
\addcontentsline{toc}{subsubsection}{14. Desarrollo estimado}
Separar:

\begin{itemize}
\item platform;
\item sensor payload;
\item NDE;
\item software;
\item qualification;
\item field validation.
\end{itemize}

Dar rango y justificar.


\subsubsection*{15. Business model}
\addcontentsline{toc}{subsubsection}{15. Business model}
Comparar:

\begin{itemize}
\item venta de robot;
\item lease;
\item service;
\item per-inspection;
\item annual contract;
\item data/analytics subscription.
\end{itemize}


\subsubsection*{16. Cliente y comprador}
\addcontentsline{toc}{subsubsection}{16. Cliente y comprador}
\begin{itemize}
\item utility;
\item OEM;
\item outage contractor;
\item NDE company;
\item insurance;
\item asset owner.
\end{itemize}

Identificar quién firma el PO.


\subsubsection*{17. Frecuencia de compra/inspección}
\addcontentsline{toc}{subsubsection}{17. Frecuencia de compra/inspección}
Fundamentar.


\subsubsection*{18. Market size}
\addcontentsline{toc}{subsubsection}{18. Market size}
Calcular solamente si hay datos suficientes.

Separar:

\begin{itemize}
\item global;
\item USA;
\item Latinoamérica;
\item Argentina.
\end{itemize}


\subsubsection*{19. Precio/ticket}
\addcontentsline{toc}{subsubsection}{19. Precio/ticket}
Usar:

\begin{itemize}
\item precios publicados;
\item RFP;
\item contratos;
\item quotes;
\item comparables.
\end{itemize}

No inventar.


\subsubsection*{20. Reutilización de nuestra plataforma RI}
\addcontentsline{toc}{subsubsection}{20. Reutilización de nuestra plataforma RI}
Puntuar:

\begin{itemize}
\item locomoción;
\item Kria/compute;
\item cámaras;
\item iluminación;
\item tether;
\item FPGA;
\item encoders;
\item software;
\item mapping;
\item reporting;
\item safety architecture.
\end{itemize}


\subsubsection*{21. Moat posible}
\addcontentsline{toc}{subsubsection}{21. Moat posible}
\begin{itemize}
\item mecánica;
\item NDE;
\item dataset;
\item calibration;
\item software;
\item qualification;
\item workflow/service;
\item patents.
\end{itemize}


\subsubsection*{22. Go / No-Go}
\addcontentsline{toc}{subsubsection}{22. Go / No-Go}
Conclusión breve y contundente.


\subsubsection*{23. Confianza}
\addcontentsline{toc}{subsubsection}{23. Confianza}
\begin{itemize}
\item Alta;
\item Media;
\item Baja.
\end{itemize}

Explicar qué dato falta.


\section*{11. Tamaño de mercado: metodología obligatoria}
\addcontentsline{toc}{section}{11. Tamaño de mercado: metodología obligatoria}
No usar como base reportes genéricos de ``robot inspection market''.

Construir \textbf{bottom-up}.


\subsection*{Fuentes preferidas}
\addcontentsline{toc}{subsection}{Fuentes preferidas}

\subsubsection*{USA}
\addcontentsline{toc}{subsubsection}{USA}
\begin{itemize}
\item EIA;
\item NRC;
\item DOE;
\item FERC;
\item NERC;
\item utilities;
\item EPRI.
\end{itemize}


\subsubsection*{Nuclear global}
\addcontentsline{toc}{subsubsection}{Nuclear global}
\begin{itemize}
\item IAEA PRIS;
\item WNA sólo como secundaria.
\end{itemize}


\subsubsection*{Argentina}
\addcontentsline{toc}{subsubsection}{Argentina}
\begin{itemize}
\item CAMMESA;
\item Secretaría de Energía;
\item ENRE;
\item Nucleoeléctrica;
\item Transener;
\item empresas generadoras;
\item operadores de ciclo combinado/hidro.
\end{itemize}


\subsubsection*{Latinoamérica}
\addcontentsline{toc}{subsubsection}{Latinoamérica}
\begin{itemize}
\item ministerios/ISO/operadores eléctricos;
\item utilities;
\item OLADE;
\item CIER;
\item bases públicas verificables.
\end{itemize}


\subsection*{Cálculo sugerido}
\addcontentsline{toc}{subsection}{Cálculo sugerido}
Por oportunidad:

\texttt{Activos addressables $\times$ inspecciones/año $\times$ ticket de servicio = TAM de servicios}

Mostrar:

\begin{itemize}
\item low case;
\item base case;
\item high case.
\end{itemize}

Si falta información crítica, NO fabricar el TAM.


\section*{12. Economía del cliente}
\addcontentsline{toc}{section}{12. Economía del cliente}
El análisis debe mirar la decisión desde el cliente.

Construir cuando sea posible:

\texttt{Valor creado = costo convencional evitado + outage evitado + riesgo evitado + cobertura adicional - costo de robotización}

Separar los componentes.

No asignar valor monetario arbitrario a seguridad.

Buscar ejemplos reales de:

\begin{itemize}
\item días de outage;
\item scaffolding;
\item excavation;
\item dewatering;
\item tank draining;
\item rotor pull;
\item human confined-space entry;
\item dive operation;
\item heavy-lift operation.
\end{itemize}


\section*{13. Ranking: NO usar una puntuación arbitraria sin explicación}
\addcontentsline{toc}{section}{13. Ranking: NO usar una puntuación arbitraria sin explicación}
Crear al menos dos rankings.


\subsection*{Ranking A — mejor próximo producto / bajo riesgo}
\addcontentsline{toc}{subsection}{Ranking A — mejor próximo producto / bajo riesgo}
Pesos iniciales sugeridos:

\begin{itemize}
\item dolor económico / avoided cost: 20\%
\item competencia / whitespace: 15\%
\item reutilización RI: 15\%
\item factibilidad técnica: 15\%
\item ticket/margen: 10\%
\item frecuencia/recurrencia: 10\%
\item tamaño de mercado: 10\%
\item overlap de clientes/canal comercial: 5\%
\end{itemize}


\subsection*{Ranking B — mayor moat / upside a largo plazo}
\addcontentsline{toc}{subsection}{Ranking B — mayor moat / upside a largo plazo}
\begin{itemize}
\item whitespace/IP: 20\%
\item pain económico: 20\%
\item barrera tecnológica defendible: 20\%
\item mercado: 15\%
\item ticket: 10\%
\item recurrencia: 5\%
\item reutilización RI: 5\%
\item regulatory burden inverso: 5\%
\end{itemize}

Claude puede cambiar los pesos si lo justifica.


\subsubsection*{Importante}
\addcontentsline{toc}{subsubsection}{Importante}
Mostrar una \textbf{sensitivity analysis}:
¿cambia el ganador si se modifica $\pm$20\% el peso de competencia, desarrollo o mercado?

Los scores son \textbf{síntesis de ingeniería/comercial}, no datos medidos. Etiquetarlos como tales.


\section*{14. Gráficos requeridos}
\addcontentsline{toc}{section}{14. Gráficos requeridos}
Crear gráficos sólo si agregan información.

Recomendados:

\begin{enumerate}
\item \textbf{Value vs market maturity}
\begin{itemize}
\item X: madurez/competencia
\item Y: valor económico para cliente
\item tamaño de burbuja: mercado estimado (sólo si es defendible)
\end{itemize}
\item \textbf{Reuse from RI vs development difficulty}
\item \textbf{Opportunity score comparison}
\item \textbf{Timeline of prior art / commercialization}
Para los top 3.
\item \textbf{Customer economics}
Convencional vs robotic cuando existan números reales.
\item \textbf{Market landscape}
Número/cluster de proveedores por oportunidad.
\end{enumerate}

Toda gráfica debe:

\begin{itemize}
\item indicar fuentes;
\item diferenciar datos de scores;
\item no inventar precisión numérica.
\end{itemize}


\section*{15. Imágenes en el HTML}
\addcontentsline{toc}{section}{15. Imágenes en el HTML}
Usar imágenes cuando ayuden a entender:

\begin{itemize}
\item geometría del activo;
\item robot existente;
\item mecanismo de acceso;
\item herramienta NDE;
\item arquitectura propuesta.
\end{itemize}


\subsection*{Preferencia}
\addcontentsline{toc}{subsection}{Preferencia}
\begin{enumerate}
\item EPRI;
\item OEM/fabricante;
\item utility/case study;
\item paper;
\item patente.
\end{enumerate}

Evitar stock photos decorativas.

Cada imagen debe tener:

\begin{itemize}
\item caption;
\item fabricante/sistema;
\item fuente;
\item URL;
\item fecha de consulta.
\end{itemize}

Si copyright/licencia no permite embebido confiable:

\begin{itemize}
\item no copiar;
\item usar enlace o crear un esquema original claramente marcado como \textbf{``Esquema del autor''}.
\end{itemize}

No hotlinkear imágenes frágiles si el HTML debe quedar autocontenido. Si es legal/técnicamente apropiado, guardar assets locales en \texttt{/assets}.


\section*{16. Prior-art y patentes}
\addcontentsline{toc}{section}{16. Prior-art y patentes}
Para los \textbf{Top 5} hacer búsqueda de patentes formal.

Buscar:

\begin{itemize}
\item Google Patents;
\item Espacenet;
\item WIPO Patentscope;
\item USPTO.
\end{itemize}

Keywords combinadas:

\begin{itemize}
\item asset + robotic inspection
\item asset + crawler
\item asset + remote inspection
\item asset + NDE
\item asset + wall climbing
\item asset + snake robot
\item asset + inspection vehicle
\item asset + manipulator
\end{itemize}

Registrar:

\begingroup\footnotesize
\setlength{\tabcolsep}{4pt}
\setlength{\emergencystretch}{3em}
\hyphenpenalty=50 \exhyphenpenalty=50
\setlength{\colw}{\dimexpr(\linewidth-12\tabcolsep)/6\relax}
\begin{longtable}{@{}L L L L L L@{}}
\toprule
\textbf{Familia} & \textbf{Prioridad} & \textbf{Titular} & \textbf{Qué reivindica} & \textbf{Estado} & \textbf{Riesgo para nosotros} \\
\midrule
\endfirsthead
\toprule
\textbf{Familia} & \textbf{Prioridad} & \textbf{Titular} & \textbf{Qué reivindica} & \textbf{Estado} & \textbf{Riesgo para nosotros} \\
\midrule
\endhead
\bottomrule
\endfoot
\multicolumn{6}{@{}l@{}}{\itshape Tabla a completar durante la investigación.} \\
\end{longtable}
\endgroup

No hacer una opinión legal de FTO; llamarlo \textbf{``screening preliminar de IP''}.


\section*{17. Investigación de competidores}
\addcontentsline{toc}{section}{17. Investigación de competidores}
Para cada top opportunity buscar exhaustivamente:


\subsubsection*{OEM de robot}
\addcontentsline{toc}{subsubsection}{OEM de robot}
¿Quién fabrica el hardware?


\subsubsection*{NDE vendor}
\addcontentsline{toc}{subsubsection}{NDE vendor}
¿Quién fabrica/posee el sensor?


\subsubsection*{Service company}
\addcontentsline{toc}{subsubsection}{Service company}
¿Quién realmente vende la inspección?


\subsubsection*{Integrator}
\addcontentsline{toc}{subsubsection}{Integrator}
¿Quién combina ambas cosas?


\subsubsection*{Utility-developed technology}
\addcontentsline{toc}{subsubsection}{Utility-developed technology}
¿La utility desarrolló internamente una solución?


\subsubsection*{University spin-off}
\addcontentsline{toc}{subsubsection}{University spin-off}
¿Existe un spin-off comercial?


\subsubsection*{Acquisition history}
\addcontentsline{toc}{subsubsection}{Acquisition history}
¿El producto cambió de dueño/nombre?


\subsubsection*{Current status}
\addcontentsline{toc}{subsubsection}{Current status}
A fecha 2026:

\begin{itemize}
\item active;
\item discontinued;
\item acquired;
\item no longer marketed;
\item unclear.
\end{itemize}

Buscar además:

\begin{itemize}
\item LinkedIn sólo como lead;
\item Wayback / archived brochures si un producto desapareció;
\item patents;
\item trade show proceedings;
\item RFP award documents.
\end{itemize}


\section*{18. RFPs y contratos: fuente de verdad comercial}
\addcontentsline{toc}{section}{18. RFPs y contratos: fuente de verdad comercial}
Buscar RFP/RFQ/bids/awards reales para cada top opportunity.

Términos:

\begin{itemize}
\item ``inspection services''
\item ``robotic inspection''
\item ``isophase bus''
\item ``HRSG inspection''
\item ``vault inspection''
\item ``buried piping''
\item ``penstock inspection''
\item ``dry cask inspection''
\item ``crawler''
\item ``remote visual inspection''
\end{itemize}

Preferir portales oficiales:

\begin{itemize}
\item utility procurement;
\item state procurement;
\item federal procurement;
\item municipal utilities;
\item nuclear procurement si es público.
\end{itemize}

Extraer:

\begin{itemize}
\item scope;
\item frecuencia;
\item requisitos;
\item duración;
\item award amount si aparece;
\item vendor ganador.
\end{itemize}

Esto es mucho más útil que un ``market size report''.


\section*{19. Argentina y Latinoamérica}
\addcontentsline{toc}{section}{19. Argentina y Latinoamérica}
Para los \textbf{Top 3} incluir capítulo específico.

Preguntas:

\begin{enumerate}
\item ¿Cuántos activos existen en Argentina?
\item ¿Quiénes los operan?
\item ¿Quién hoy hace la inspección?
\item ¿Contratan OEM internacional o servicio local?
\item ¿El robot podría prestar servicio regional desde Argentina?
\item ¿Qué importación/calibración requiere?
\item ¿Hay integradores/NDE partners locales?
\item ¿Hay centrales con equipos GE/Siemens/Mitsubishi/ABB donde el canal comercial sea similar al RI?
\item ¿Qué mercados vecinos son naturales?
\begin{itemize}
\item Brasil
\item Chile
\item Uruguay
\item Paraguay
\item Perú
\item Colombia
\item México
\end{itemize}
\end{enumerate}

No inventar cantidades. Usar CAMMESA/operadores/utilities.


\section*{20. Investigación de costos de desarrollo}
\addcontentsline{toc}{section}{20. Investigación de costos de desarrollo}
Para cada shortlist estimar separadamente:

\begin{itemize}
\item diseño mecánico;
\item prototipo;
\item electrónica;
\item sensors/NDE;
\item tether/power;
\item software/control;
\item autonomous navigation;
\item data/reporting;
\item test rig;
\item mock-up;
\item field qualification;
\item certification;
\item tooling/manufacturing.
\end{itemize}

Rangos amplios son preferibles a falsa precisión.

Diferenciar:

\begin{itemize}
\item costo de construir 1 prototipo;
\item NRE para industrialización;
\item costo unitario aproximado;
\item costo de field service kit.
\end{itemize}

Buscar precios reales de componentes críticos cuando sea posible.


\section*{21. Estructura del informe HTML final}
\addcontentsline{toc}{section}{21. Estructura del informe HTML final}

\subsection*{Portada}
\addcontentsline{toc}{subsection}{Portada}
\textbf{Oportunidades de Robótica de Inspección Industrial para Generación y Utilities} \\
\emph{State of the art, landscape competitivo y evaluación de negocio — 2026}


\subsection*{Executive Summary}
\addcontentsline{toc}{subsection}{Executive Summary}
Máximo \textasciitilde{}3 páginas:

\begin{itemize}
\item top 3;
\item por qué;
\item qué NO desarrollar;
\item biggest unknowns;
\item decisión recomendada.
\end{itemize}


\subsection*{1. Objetivo y alcance}
\addcontentsline{toc}{subsection}{1. Objetivo y alcance}

\subsection*{2. Metodología y evidencia}
\addcontentsline{toc}{subsection}{2. Metodología y evidencia}
\begin{itemize}
\item jerarquía;
\item search coverage;
\item limitaciones.
\end{itemize}


\subsection*{3. Qué hace atractivo un negocio de inspection robotics}
\addcontentsline{toc}{subsection}{3. Qué hace atractivo un negocio de inspection robotics}
\begin{itemize}
\item economics;
\item outage;
\item service model;
\item FME;
\item qualification.
\end{itemize}


\subsection*{4. Mapa completo de oportunidades}
\addcontentsline{toc}{subsection}{4. Mapa completo de oportunidades}
Taxonomía por activo/sector.


\subsection*{5. Candidate cards}
\addcontentsline{toc}{subsection}{5. Candidate cards}
Una sección completa por robot.


\subsection*{6. Competitive landscape}
\addcontentsline{toc}{subsection}{6. Competitive landscape}
Tabla transversal de vendors/products.


\subsection*{7. Patent / IP landscape}
\addcontentsline{toc}{subsection}{7. Patent / IP landscape}

\subsection*{8. Mercado}
\addcontentsline{toc}{subsection}{8. Mercado}
\begin{itemize}
\item global;
\item USA;
\item LATAM;
\item Argentina.
\end{itemize}


\subsection*{9. Economics}
\addcontentsline{toc}{subsection}{9. Economics}
\begin{itemize}
\item customer ROI;
\item ticket;
\item service model.
\end{itemize}


\subsection*{10. Comparison matrix}
\addcontentsline{toc}{subsection}{10. Comparison matrix}

\subsection*{11. Rankings}
\addcontentsline{toc}{subsection}{11. Rankings}
\begin{itemize}
\item fast-follow;
\item deep-tech/moat;
\item sensitivity.
\end{itemize}


\subsection*{12. Top 3 — Concept briefs}
\addcontentsline{toc}{subsection}{12. Top 3 — Concept briefs}
Para cada uno:

\begin{itemize}
\item target requirements;
\item MVP;
\item differentiation;
\item validation plan;
\item 12–24 month roadmap;
\item customer discovery.
\end{itemize}


\subsection*{13. Do-not-pursue-now}
\addcontentsline{toc}{subsection}{13. Do-not-pursue-now}
Explicar por qué.


\subsection*{14. Unknowns / due diligence}
\addcontentsline{toc}{subsection}{14. Unknowns / due diligence}
Lista accionable.


\subsection*{15. Conclusiones}
\addcontentsline{toc}{subsection}{15. Conclusiones}

\subsection*{Anexo A — Search log}
\addcontentsline{toc}{subsection}{Anexo A — Search log}
Keywords, bases y fecha.


\subsection*{Anexo B — Competitor register}
\addcontentsline{toc}{subsection}{Anexo B — Competitor register}

\subsection*{Anexo C — Patent register}
\addcontentsline{toc}{subsection}{Anexo C — Patent register}

\subsection*{Anexo D — Source register}
\addcontentsline{toc}{subsection}{Anexo D — Source register}

\section*{22. Formato de referencias}
\addcontentsline{toc}{section}{22. Formato de referencias}
Usar referencias junto a la información.

Ejemplo:

\begin{callout}
EPRI reporta que una utility pagó aproximadamente USD 100.000 por una implementación de inspección robótica de isophase bus y que la ejecución tomó menos de un día [EPRI-MTA-029].
\end{callout}

Al hacer click o hover debería ser fácil localizar la referencia.

Al final:

\textbf{[EPRI-MTA-029]} Electric Power Research Institute, \emph{Reduce Inspection Cost for Isophase Bus Duct Using Robotic Crawler}, Plant Modernization Toolbox, consultado 2026-08-31. URL...

No poner una bibliografía desconectada donde resulte imposible saber qué afirmación sostiene cada fuente.


\subsubsection*{Cada referencia debe registrar}
\addcontentsline{toc}{subsubsection}{Cada referencia debe registrar}
\begin{itemize}
\item autor/organización;
\item título;
\item fecha;
\item Product ID/DOI/patent;
\item URL;
\item fecha de acceso;
\item evidence level.
\end{itemize}


\section*{23. Paquete de evidencia inicial verificado}
\addcontentsline{toc}{section}{23. Paquete de evidencia inicial verificado}
Estas fuentes son \textbf{seeds}, no reemplazan tu investigación.


\subsection*{EPRI — oportunidad y landscape}
\addcontentsline{toc}{subsection}{EPRI — oportunidad y landscape}
\begin{enumerate}
\item EPRI 3002023899 — Robotic Process Automation for Nuclear Power Plants \\
\url{https://restservice.epri.com/publicdownload/000000003002023899/0/Product}
\item EPRI Unmanned Mobile Technologies — Product Index \\
\url{https://transmission.epri.com/p37_substations/robotics/mobiletechgrp/products/}
\item EPRI underground transmission research updates \\
\url{https://transmission.epri.com/p36_underground/public/p36001_design_construction/research_updates/}
\item EPRI underground transmission applications \\
\url{https://transmission.epri.com/p36_underground/public/p36_applications/}
\item EPRI underground transmission contacts \\
\url{https://transmission.epri.com/p36_underground/leads/}
\end{enumerate}


\subsection*{Isophase bus}
\addcontentsline{toc}{subsection}{Isophase bus}
\begin{enumerate}[start=6]
\item EPRI MTA-MA-029 \\
\url{https://nuclearplantmod.epri.com/MTA-MA-029}
\item EPRI 1015057 — Isolated Phase Bus Maintenance Guide \\
\url{https://restservice.epri.com/publicdownload/000000000001015057/0/Product}
\item Trident NGS \\
\url{https://www.tridentngs.com/}
\end{enumerate}


\subsection*{HRSG}
\addcontentsline{toc}{subsection}{HRSG}
\begin{enumerate}[start=9]
\item EPRI 1017635 — Study for Snake Robot Technology for Inspection of Headers and Tubes in HRSGs \\
\url{https://restservice.epri.com/publicdownload/000000000001017635/0/Product}
\item TesTex HRSG \\
\url{https://testex-ndt.com/services/hrsg/}
\item TesTex IAT \\
\url{https://testex-ndt.com/products/iat-internal-access-tool-for-hrsg-inspections/}
\item TesTex Claw \\
\url{https://testex-ndt.com/products/claw-bfet-inspection-tool/}
\end{enumerate}


\subsection*{Buried piping}
\addcontentsline{toc}{subsection}{Buried piping}
\begin{enumerate}[start=13]
\item EPRI MTA-MA-017 \\
\url{https://nuclearplantmod.epri.com/MTA-MA-017}
\item EPRI 1025272 \\
\url{https://restservice.epri.com/publicdownload/000000000001025272/0/Product}
\item Diakont — first buried-pipe inline inspection case \\
\url{https://diakont.com/case-studies/nuclear-solutions/first-buried-piping-in-line-inspection/}
\end{enumerate}


\subsection*{Hydro}
\addcontentsline{toc}{subsection}{Hydro}
\begin{enumerate}[start=16]
\item EPRI 1007576 / TR-113584-V7 — ROV Technology at Hydro Facilities \\
\url{https://restservice.epri.com/publicdownload/000000000001007576/0/Product}
\end{enumerate}


\subsection*{Dry cask}
\addcontentsline{toc}{subsection}{Dry cask}
\begin{enumerate}[start=17]
\item EPRI Journal — wall-climbing dry cask robots \\
\url{https://eprijournal.com/wall-climbing-robots-inspect-nuclear-storage-casks/}
\item Robotic Technologies of Tennessee \\
\url{https://www.robotictechtn.com/}
\end{enumerate}


\subsection*{Emerging}
\addcontentsline{toc}{subsection}{Emerging}
\begin{enumerate}[start=19]
\item EPRI Generation Reports \\
\url{https://dx-wiki.epri.com/Generation_Reports}
\item EPRI Tech Portal \\
\url{https://techportal.epri.com/}
\end{enumerate}


\section*{24. Hallazgos seed que NO deben perderse}
\addcontentsline{toc}{section}{24. Hallazgos seed que NO deben perderse}
Estos hallazgos ya fueron verificados en fuentes primarias/actuales, pero debés volver a leer el contexto antes de utilizarlos.


\subsection*{IPB}
\addcontentsline{toc}{subsection}{IPB}
EPRI MTA-MA-029 informa:

\begin{itemize}
\item vendor-supported inspection \textasciitilde{}USD 100k en experiencia de una utility;
\item \textless{}1 día de ejecución;
\item reducción de costo de aproximadamente 10$\times$ frente a método manual en ese caso;
\item payback inmediato;
\item commercial readiness alta.
\end{itemize}

Esto demuestra \textbf{mercado y valor}, pero también que un ``camera crawler'' solo ya existe.


\subsection*{Buried piping}
\addcontentsline{toc}{subsection}{Buried piping}
EPRI MTA-MA-017 informa:

\begin{itemize}
\item ahorro esperado \textgreater{}USD 1 M y \textless{}USD 5 M por uso;
\item implementación \textless{}USD 1 M;
\item payback \textless{}1 año/inmediato;
\item ya comercial en nuclear.
\end{itemize}

Esto indica \textbf{gran valor pero madurez alta}.


\subsection*{HRSG}
\addcontentsline{toc}{subsection}{HRSG}
EPRI 1017635 confirma que:

\begin{itemize}
\item ciertos tubos interiores no pueden inspeccionarse fácilmente con NDE convencional sin acceso destructivo;
\item snake robotics fue estudiada para ese problema.
\end{itemize}

Pero TesTex hoy ofrece:

\begin{itemize}
\item IAT;
\item RFET + video;
\item crawler en header;
\item múltiples work sites.
\end{itemize}

Por lo tanto, la oportunidad HRSG debe reevaluarse contra un incumbente real.


\subsection*{Underground transmission vault}
\addcontentsline{toc}{subsection}{Underground transmission vault}
EPRI mantiene el proyecto activo en 2025–2027:

\begin{itemize}
\item prototipos;
\item lab;
\item utility demos;
\item objetivo de seguridad y reducción de outage.
\end{itemize}

Esto es una señal de \textbf{problema vigente y tecnología aún en evolución}.


\section*{25. Preguntas que el informe DEBE responder al final}
\addcontentsline{toc}{section}{25. Preguntas que el informe DEBE responder al final}
\begin{enumerate}
\item \textbf{Si tuviéramos que invertir hoy en un solo robot después del RI, ¿cuál sería?}
\item ¿Por qué un cliente pagaría por él?
\item ¿Qué costo concreto le evitamos?
\item ¿Qué ticket de servicio puede sostener?
\item ¿Cuántas oportunidades de servicio existen por año?
\item ¿Quién ya hace esto?
\item ¿Por qué nuestro producto sería mejor/diferente?
\item ¿Qué tan difícil es copiarlo?
\item ¿Qué IP puede bloquearlo?
\item ¿Qué parte del RI actual podemos reutilizar?
\item ¿Cuánto desarrollo falta?
\item ¿Cuál es el primer MVP que un cliente pagaría por probar?
\item ¿Qué validación técnica sería suficiente?
\item ¿Cuál es el principal mecanismo de fracaso?
\item ¿Qué robot NO deberíamos desarrollar aunque parezca atractivo?
\item ¿Qué oportunidad nueva apareció durante la investigación que no estaba en este brief?
\item ¿Qué información todavía no existe públicamente y requiere entrevistas?
\item ¿A qué 10 personas/organizaciones deberíamos contactar primero?
\end{enumerate}


\section*{26. Customer discovery recomendado}
\addcontentsline{toc}{section}{26. Customer discovery recomendado}
Para los Top 3 generar una lista concreta de entrevistas.

Tipos de persona:

\begin{itemize}
\item generator maintenance manager;
\item outage manager;
\item NDE manager;
\item plant manager;
\item transmission cable engineer;
\item HRSG engineer;
\item OEM service manager;
\item EPRI technical lead;
\item NDE service provider;
\item insurance/risk engineer cuando aplique.
\end{itemize}

Crear 8–12 preguntas por segmento orientadas a:

\begin{itemize}
\item costo actual;
\item frecuencia;
\item pain;
\item fallas;
\item willingness to pay;
\item acceptance requirements;
\item tool recovery/FME;
\item procurement;
\item incumbent.
\end{itemize}


\section*{27. Requisitos visuales del HTML}
\addcontentsline{toc}{section}{27. Requisitos visuales del HTML}
El informe anterior funcionó bien visualmente; mantener un nivel profesional.


\subsubsection*{Estilo}
\addcontentsline{toc}{subsubsection}{Estilo}
\begin{itemize}
\item limpio;
\item técnico;
\item ejecutivo;
\item legible;
\item responsive;
\item tablas con sticky header cuando sean largas;
\item TOC lateral o superior;
\item callouts de ``Hallazgo'', ``Riesgo'', ``Inferencia'', ``Dato EPRI'';
\item no saturar de colores;
\item gráficos vectoriales/SVG cuando convenga.
\end{itemize}


\subsubsection*{Componentes útiles}
\addcontentsline{toc}{subsubsection}{Componentes útiles}
\begin{itemize}
\item executive scorecards;
\item comparison matrix;
\item expandable source notes;
\item evidence badges A/B/C/D;
\item maturity badges M0–M5;
\item confidence badges.
\end{itemize}


\subsubsection*{Evitar}
\addcontentsline{toc}{subsubsection}{Evitar}
\begin{itemize}
\item marketing genérico;
\item frases sin sustento;
\item imágenes decorativas;
\item gráficas 3D;
\item tablas ilegibles;
\item ``AI generated'' filler.
\end{itemize}


\section*{28. Archivos finales a entregar}
\addcontentsline{toc}{section}{28. Archivos finales a entregar}
Crear:

\begin{enumerate}
\item \texttt{robot\_opportunities\_2026.html}
\begin{itemize}
\item informe completo en español.
\end{itemize}
\item \texttt{robot\_opportunities\_sources.md}
\begin{itemize}
\item registro de fuentes con notas de qué respalda cada una.
\end{itemize}
\item \texttt{robot\_opportunities\_evidence.csv}
Columnas sugeridas:
\begin{itemize}
\item candidate
\item claim
\item value
\item unit
\item source\_id
\item source\_type
\item evidence\_level
\item url
\item access\_date
\item confidence
\item notes
\end{itemize}
\item \texttt{robot\_opportunities\_competitors.csv}
\begin{itemize}
\item producto;
\item empresa;
\item aplicación;
\item status;
\item país;
\item tech;
\item dimensions;
\item payload;
\item deployments;
\item price;
\item source.
\end{itemize}
\item \texttt{robot\_opportunities\_patents.csv}
\item \texttt{/assets}
\begin{itemize}
\item imágenes permitidas/locales;
\item gráficos generados;
\item diagramas.
\end{itemize}
\item \texttt{research\_log.md}
\begin{itemize}
\item consultas realizadas;
\item bases consultadas;
\item leads descartados;
\item fuentes que no pudieron obtenerse.
\end{itemize}
\end{enumerate}


\section*{29. Criterio de finalización}
\addcontentsline{toc}{section}{29. Criterio de finalización}
NO dar por terminado el informe porque ``hay suficiente información''.

Antes de cerrar, ejecutar este checklist:

\begin{itemize}[label=$\square$,leftmargin=1.6em]
\item Se recorrió EPRI Product Index completo de robotics/mobile technologies.
\item Se recorrió EPRI Plant Modernization Toolbox buscando robotics/inspection.
\item Se buscó cada activo + robot en EPRI.
\item Se buscó cada activo sin la palabra robot para encontrar métodos competidores.
\item Se investigaron al menos los Top 7 candidatos.
\item Se buscaron candidatos nuevos.
\item Se hizo competitor search profundo de Top 5.
\item Se hizo patent screening de Top 5.
\item Se buscaron RFPs/contratos de Top 3.
\item Se verificó status 2026 de cada competidor principal.
\item Se buscó precio/ticket cuando exista.
\item Cada número importante tiene referencia inmediata.
\item Cada claim de fabricante está etiquetado.
\item Cada inferencia está etiquetada.
\item Se documentó evidencia negativa.
\item Se hizo market sizing bottom-up o se explicó por qué no es defendible.
\item Se hizo capítulo Argentina/LATAM para Top 3.
\item Se hizo sensitivity analysis del ranking.
\item Se definieron Top 3.
\item Se definieron ``Do not pursue now''.
\item Se listaron unknowns.
\item Se listaron entrevistas/contactos prioritarios.
\item Se revisó que el HTML sea autocontenido y abra correctamente.
\end{itemize}


\section*{30. Instrucción final}
\addcontentsline{toc}{section}{30. Instrucción final}
\textbf{Buscá más de lo que te estoy dando acá.}

Este brief contiene hipótesis y fuentes seed, no las conclusiones.

Si durante la investigación encontrás que:

\begin{itemize}
\item IPB está saturado;
\item TesTex cerró el gap de HRSG;
\item vault robotics no tiene economics;
\item existe un robot desconocido que invalida nuestra propuesta;
\item otra aplicación EPRI ofrece un negocio mucho mejor;
\end{itemize}

\textbf{cambiá el ranking.}

La prioridad es no enamorarse de una solución.

El resultado debe permitir tomar una decisión de inversión con tres cosas claramente separadas:

\begin{enumerate}
\item \textbf{lo que sabemos;}
\item \textbf{lo que inferimos;}
\item \textbf{lo que todavía debemos validar.}
\end{enumerate}

Y la conclusión debe responder sin ambigüedad:

\begin{callout}
\textbf{¿Cuál debería ser el próximo robot que desarrollemos, y por qué?}
\end{callout}


\end{document}
